% archon:physics

\chapter{Finite Hamiltonian rescaling}
\label{ch:ArchonPhysics_HamiltonianScaling}

This clean-room module formalizes only the finite-dimensional algebra of the
periodic polynomial Hamiltonian.  The configuration carrier is the lattice
configuration supplied by Chapter~\ref{ch:ArchonPhysics_Lattice}; position and
momentum are its named scalar coordinate projections in fixed units.  The
energy-density parameter $\epsilon$ is positive and dimensionless after this
fixed choice of energy unit.  No statement below asserts a kinetic equation,
thermalization, a global Hamiltonian flow, or an equilibration time law.

\begin{definition}\leanok
[Hamilton-equation predicate]
\label{def:hamiltonian:satisfies-equations}
\lean{ArchonPhysics.HamiltonianScaling.SatisfiesHamiltonEquations}
For a Hamiltonian $H(t,p,q)$ and momentum and position paths $p,q$, this
predicate means that the Physlib Hamilton-equation residual is zero.
\end{definition}
\begin{proof}

\leanok
This is a named wrapper around the supplied classical-mechanics residual, so
that the API records its physical role rather than treating the condition as an
unexplained scalar equation.
\end{proof}

\begin{theorem}\leanok
[Hamilton-equation residual characterization]
\label{thm:hamiltonian:satisfies-iff}
\lean{ArchonPhysics.HamiltonianScaling.satisfiesHamiltonEquations_iff}
\uses{def:hamiltonian:satisfies-equations}
The named predicate holds exactly when
$\operatorname{hamiltonEqOp}(H,p,q)=0$.
\end{theorem}
\begin{proof}

\leanok
Unfold the wrapper defining the predicate.  Both sides are the same residual
equation.
\end{proof}

\begin{definition}\leanok
[Configuration rescaling]
\label{def:hamiltonian:rescale-configuration}
\lean{ArchonPhysics.HamiltonianScaling.rescaleConfiguration}
\uses{def:lattice:configuration}
For an energy-density scale $\epsilon$, rescale every coordinate of a lattice
configuration by $\sqrt{\epsilon}$.
\end{definition}
\begin{proof}

\leanok
The rescaling is the pointwise map $q_i\mapsto\sqrt{\epsilon}\,q_i$.
\end{proof}

\begin{definition}\leanok
[Effective nonlinear coupling]
\label{def:hamiltonian:effective-coupling}
\lean{ArchonPhysics.HamiltonianScaling.effectiveCoupling}
For polynomial order $n$, define
$g(\lambda,\epsilon,n)=\lambda\epsilon^{((n: \mathbb R)-2)/2}$.
\end{definition}
\begin{proof}

\leanok
This is the scalar coupling induced by extracting one factor of $\epsilon$
from the degree-$n$ interaction energy.
\end{proof}

\begin{definition}\leanok
[Finite periodic Hamiltonian]
\label{def:hamiltonian:lattice-hamiltonian}
\lean{ArchonPhysics.HamiltonianScaling.latticeHamiltonian}
\uses{def:lattice:mass, def:lattice:forward-difference}
For a positive mass profile $m$, set
\[
 H_{m,n,\lambda}(p,q)=\sum_i\left(
 \frac{p_i^2}{2m_i}+\frac{(Dq)_i^2}{2}
 +\frac{\lambda}{n}(Dq)_i^n\right).
\]
The first term is kinetic energy, the second is the harmonic bond energy, and
the final term is the degree-$n$ nearest-neighbour interaction energy.
\end{definition}
\begin{proof}

\leanok
The locked Hamiltonian is this finite sum over cyclic sites; it uses the
positive mass field and the periodic forward-difference operator from the
lattice API.
\end{proof}

\begin{theorem}\leanok
[Rescaling data specification]
\label{thm:hamiltonian:rescaling-data}
\lean{ArchonPhysics.HamiltonianScaling.rescalingData_spec}
\uses{def:hamiltonian:lattice-hamiltonian,
      def:hamiltonian:rescale-configuration,
      def:hamiltonian:effective-coupling}
The Hamiltonian has the displayed finite-sum formula, configuration rescaling
is pointwise multiplication by $\sqrt\epsilon$, and the effective coupling is
$\lambda\epsilon^{((n:\mathbb R)-2)/2}$.
\end{theorem}
\begin{proof}

\leanok
Each of the three conjuncts unfolds one of the locked definitions.  The
Hamiltonian formula uses the periodic forward difference from the preceding
definition.
\end{proof}

\begin{theorem}\leanok
[Finite Hamiltonian rescaling]
\label{thm:hamiltonian:lattice-rescale}
\lean{ArchonPhysics.HamiltonianScaling.latticeHamiltonian_rescale}
\uses{thm:hamiltonian:rescaling-data}
For $n\ge3$ and $\epsilon>0$,
\[
 H_{m,n,\lambda}(\sqrt\epsilon\,p,\sqrt\epsilon\,q)
 =\epsilon\,H_{m,n,g(\lambda,\epsilon,n)}(p,q).
\]
\end{theorem}
\begin{proof}

\leanok
Expand the finite sum.  Pointwise linearity of the forward difference gives
$D(\sqrt\epsilon q)=\sqrt\epsilon Dq$.  The kinetic and harmonic squares
each contribute $\epsilon$.  For the interaction, positivity of $\epsilon$
identifies $(\sqrt\epsilon)^n$ with
$\epsilon\epsilon^{(n-2)/2}$, so its remaining coefficient is exactly $g$.
Sum these termwise identities over the finite cyclic site set.
\end{proof}

\begin{theorem}\leanok
[Inverse-square effective-coupling identity]
\label{thm:hamiltonian:effective-coupling-inv-sq}
\lean{ArchonPhysics.HamiltonianScaling.effectiveCoupling_inv_sq}
\uses{def:hamiltonian:effective-coupling}
For $n\ge3$, $\epsilon>0$, and $\lambda\ne0$,
\[
 g(\lambda,\epsilon,n)^{-2}
 =\lambda^{-2}\epsilon^{-(n-2)}.
\]
\end{theorem}
\begin{proof}

\leanok
The positive base permits the usual real-power laws.  Invert the product
$\lambda\epsilon^{(n-2)/2}$ and square it; the two half-exponents add to
$n-2$.  The hypotheses make both inverse factors defined in the intended
nonzero sense.
\end{proof}
